\documentclass[11pt]{article}
\usepackage[a4paper,margin=2.8cm]{geometry}
\usepackage{amsmath,amssymb,amsthm,mathtools}
\usepackage{booktabs,array,enumitem}
\usepackage{microtype}
\usepackage[skip=0.75\baselineskip plus 2pt minus 1pt]{parskip}
\usepackage{titlesec}
\usepackage[font=small,labelfont=bf,margin=1em]{caption}
\usepackage[hidelinks]{hyperref}
\hypersetup{pdftitle={New bounds for small two- and four-colour Ramsey numbers},
  pdfauthor={Lachlan Dauth},
  pdfsubject={Mathematics - Combinatorics (05C55)},
  pdfkeywords={Ramsey numbers, triangle-free graphs, edge numbers e(3,k,n), quartic residue graphs, multicolour Ramsey numbers, computer-assisted proofs}}

% Layout: block paragraphs with space between them, roomier lists,
% theorems and section headings, slightly opened line spacing.
\linespread{1.06}
\setlist{itemsep=0.35em,topsep=0.5em,parsep=0pt,partopsep=0pt}
\setlist[enumerate]{leftmargin=2.2em}
\setlist[itemize]{leftmargin=1.8em}
\titlespacing*{\section}{0pt}{2.6ex plus 1ex minus .2ex}{1.4ex plus .2ex}
\titlespacing*{\subsection}{0pt}{2.0ex plus 1ex minus .2ex}{1.0ex plus .2ex}
\titlespacing*{\paragraph}{0pt}{1.2ex plus .5ex minus .2ex}{0.8em}
\renewcommand{\arraystretch}{1.15}
% abstract: block paragraphs with a gap, like the body
\renewenvironment{abstract}{\small\begin{center}\bfseries\abstractname\end{center}%
  \list{}{\setlength{\leftmargin}{2.5em}\setlength{\rightmargin}{\leftmargin}%
  \setlength{\listparindent}{0pt}\setlength{\parsep}{0.6em plus 1pt}\setlength{\topsep}{0pt}}\item\relax}{\endlist}
\newtheoremstyle{spacedplain}{1.1em}{1.1em}{\itshape}{}{\bfseries}{.}{0.6em}{}
\newtheoremstyle{spaceddef}{1.1em}{1.1em}{\normalfont}{}{\bfseries}{.}{0.6em}{}
\newtheoremstyle{spacedremark}{1.0em}{1.0em}{\normalfont}{}{\itshape}{.}{0.6em}{}
\theoremstyle{spacedplain}

\newtheorem{theorem}{Theorem}[section]
\newtheorem{lemma}[theorem]{Lemma}
\newtheorem{corollary}[theorem]{Corollary}
\newtheorem{proposition}[theorem]{Proposition}
\theoremstyle{spaceddef}
\newtheorem{definition}[theorem]{Definition}
\theoremstyle{spacedremark}
\newtheorem{remark}[theorem]{Remark}

\newcommand{\R}{\mathrm{R}}
\newcommand{\bt}{\mathbf{t}}
\newcommand{\e}{\mathbf{e}}
\newcommand{\Nr}{N_{1}}
\newcommand{\Nb}{N_{2}}
\newcommand{\Ng}{N_{3}}
% Optional layout hooks for a review copy (no effect in this version)
\ifdefined\markupcopy% Layout for the printed review copy. Loaded only when \markupcopy is defined.
% Content, numbering and wording are identical to the clean draft.
\geometry{a4paper,left=2.0cm,right=6.2cm,top=2.6cm,bottom=2.6cm,headheight=14pt,headsep=0.6cm}
\usepackage{xcolor}
\usepackage{setspace}
\setstretch{1.7}
\usepackage{etoolbox}
\AtBeginEnvironment{tabular}{\setstretch{1.15}}
\AtBeginEnvironment{thebibliography}{\setstretch{1.25}}
% wide tables may not spill into the notes column
\usepackage{adjustbox}
\let\mkoldtabular\tabular
\let\mkoldendtabular\endtabular
\renewenvironment{tabular}{\begin{adjustbox}{max width=\linewidth}\mkoldtabular}{\mkoldendtabular\end{adjustbox}}
% line numbers on every line, for referencing comments
\usepackage[mathlines]{lineno}
\renewcommand\linenumberfont{\normalfont\scriptsize\sffamily\color{black!55}}
\linenumbers
% header and footer
\usepackage{lastpage}
\usepackage{fancyhdr}
\pagestyle{fancy}
\fancyhf{}
\fancyhead[L]{\scriptsize\sffamily Review copy \textperiodcentered{} draft v5 \textperiodcentered{} printed 23 October 2026}
\fancyhead[R]{\scriptsize\sffamily page \thepage\ of \pageref{LastPage}}
\renewcommand{\headrulewidth}{0.3pt}
\fancypagestyle{plain}{\fancyhf{}\fancyhead[L]{\scriptsize\sffamily Review copy \textperiodcentered{} draft v5 \textperiodcentered{} printed 23 October 2026}\fancyhead[R]{\scriptsize\sffamily page \thepage\ of \pageref{LastPage}}}
% ruled notes column on every page
\usepackage{eso-pic,tikz}
\AddToShipoutPictureBG{%
  \begin{tikzpicture}[remember picture,overlay]
    \draw[black!35,line width=0.4pt]
      ([xshift=-5.85cm,yshift=-2.6cm]current page.north east) -- ([xshift=-5.85cm,yshift=2.6cm]current page.south east);
    \node[anchor=south west,inner sep=0pt,font=\scriptsize\sffamily,text=black!50]
      at ([xshift=-5.6cm,yshift=-2.5cm]current page.north east) {Notes};
    \foreach \y in {3.4,4.2,...,26.6}
      \draw[black!22,line width=0.3pt]
        ([xshift=-5.6cm,yshift=-\y cm]current page.north east) -- ([xshift=-0.9cm,yshift=-\y cm]current page.north east);
  \end{tikzpicture}}
% a ruled page for general comments at the end
\AtEndDocument{%
  \clearpage\nolinenumbers
  \section*{General comments}
  \begin{tikzpicture}[remember picture,overlay]
    \foreach \y in {5.0,5.8,...,26.6}
      \draw[black!22,line width=0.3pt]
        ([xshift=2.0cm,yshift=-\y cm]current page.north west) -- ([xshift=-6.2cm,yshift=-\y cm]current page.north east);
  \end{tikzpicture}}
% narrower column: let paragraphs stretch rather than overrun, and break three wide displays
\setlength{\emergencystretch}{3em}
\newcommand{\mkbreak}[1]{\\}
\newcommand{\mkalign}{\\}
\fi
\providecommand{\mkbreak}[1]{#1}% this copy: keeps the given spacing; review copy: line break
\providecommand{\mkalign}{&}%     this copy: next alignment column; review copy: new row

\title{New bounds for small two- and\\ four-colour Ramsey numbers}
\author{Lachlan Dauth}
\date{23 October 2026}

\begin{document}
\maketitle

\begin{abstract}
We give upper bounds for the Ramsey numbers $\R(3,k)$ with $24\le k\le 50$, for which Radziszowski's survey lists no upper bounds.
They come from an exact bootstrap of lower bounds on the number of edges of triangle-free graphs with bounded independence number.
It starts from the published values and lower bounds of Goedgebeur and Radziszowski.
For instance, $\R(3,24)\le 201$, $\R(3,30)\le 297$ and $\R(3,50)\le 729$.
If the table of Backelin's unrefereed report is also used, the bounds are $1$ to $3$ lower for $24\le k\le 36$, except $k=35$.

Deleting vertices from the quartic residue graphs of Lindsay and Cain gives $\R(4,21)\ge 260$ and $\R(4,24)\ge 370$.
Finally, we use the four-colour form of the neighbourhood recurrence, with $\R(3,4,4)\le 73$ and $\R(4,4,4)\le 198$ from a companion paper and the survey's $\R(3,3,3,4)\le 149$.
This gives $\R(3,3,4,4)\le 442$, $\R(3,4,4,4)\le 1521$ and $\R(4,4,4,4)\le 6082$.
All certificates and witnesses are checked exactly by programs in the ancillary files.
\end{abstract}

\begin{quote}\small
\textit{2020 Mathematics Subject Classification.} Primary 05C55; Secondary 05D10, 05C35, 05C69, 05C25, 68V05.

\textit{Keywords.} Ramsey numbers; triangle-free graphs; edge numbers $e(3,k,n)$; quartic residue graphs; multicolour Ramsey numbers; computer-assisted proofs.
\end{quote}

\section{Introduction}

For integers $s,t\ge 1$, the Ramsey number $\R(s,t)$ is the least $n$ such that every red-blue colouring of the edges of $K_n$ contains a red $K_s$ or a blue $K_t$; equivalently, every graph on $n$ vertices contains a $K_s$ or an independent set of size $t$.
Multicolour numbers $\R(t_1,\dots,t_r)$ are defined in the same way.
The current state of knowledge is collected in Radziszowski's dynamic survey \cite{Rad26}, whose revision~18 we use for all comparisons, together with the later preprint \cite{Yu26}.

This paper is a companion to a study of three-colour numbers \cite{DC26}, in which local Ramsey bounds inside colour neighbourhoods are combined with double counting and with flag-algebra certificates.
The present paper gives new bounds for two-colour numbers.
The upper bounds are obtained by adapting the ladder procedure of \cite[Section~5]{DC26}, and the lower bounds by deleting vertices from quartic residue graphs.
It also derives consequences for four-colour numbers.
Our results are listed in Table~\ref{tab:results}.

\begin{table}[htbp]
\centering
\begin{tabular}{lccl}
\toprule
number & previous & new & where\\
\midrule
$\R(3,k)$, $24\le k\le 50$ & no upper bound \cite{Rad26} & Table~\ref{tab:r3k} & Theorem~\ref{thm:r3k}\\
$\R(4,21)$ & $\ge 242$ \cite{SLZL}; $\ge 257$ (\cite{Yu26} and \cite[2.3(e)]{Rad26}) & $\ge 260$ & Theorem~\ref{thm:r4k}\\
$\R(4,24)$ & $\ge 343$ (implicit in \cite{LinCa}) & $\ge 370$ & Theorem~\ref{thm:r4k}\\
$\R(3,3,4,4)$ & $\le 450$ \cite{Rad26} & $\le 442$ & Theorem~\ref{thm:four}\\
$\R(3,4,4,4)$ & $\le 1576$ \cite{Rad26} (implicit in \cite{Boz26}) & $\le 1521$ & Theorem~\ref{thm:four}\\
$\R(4,4,4,4)$ & $\le 6306$ \cite{Rad26}; $\le 6301$ (implicit in \cite{Boz26}) & $\le 6082$ & Theorem~\ref{thm:four}\\
\bottomrule
\end{tabular}
\caption{Our results.
For $\R(3,k)$ with $k\ge 24$, iterating $\R(3,k)\le\R(3,k-1)+k$, which is strict when both terms are even, or Shearer's bound \cite{Sh83}, gives weaker explicit bounds (Table~\ref{tab:r3k}).
Where the two rows of Table~\ref{tab:r3k} differ, the smaller bound also uses the unrefereed report \cite{Back3}.}
\label{tab:results}
\end{table}

The upper bounds for $\R(3,k)$ (Section~\ref{sec:r3k}) and the lower bounds for $\R(4,k)$ (Section~\ref{sec:r4k}) are computer-assisted.
Floating-point solvers and local search were used only to find certificates and witnesses, and in the negative experiments that we report.
Every certificate and witness is checked in exact arithmetic by a program that does not share code with the search (Section~\ref{sec:verify}).

The four-colour bounds (Section~\ref{sec:four}) follow by hand from $\R(3,3,3,4)\le 149$ \cite{Rad26} and from $\R(3,4,4)\le 73$ and $\R(4,4,4)\le 198$, which are proved in \cite[Theorems~8.3 and~8.2]{DC26} by certificates checked in exact arithmetic by computer.

\section{Upper bounds for \texorpdfstring{$\R(3,k)$}{R(3,k)}}\label{sec:r3k}

A $(3,k;n)$-graph is a triangle-free graph on $n$ vertices with no independent set of size $k$.
A $(3,k;n,e)$-graph is one with $e$ edges, and $e(3,k,n)$ is the least such $e$.
By analogy with \cite[Section~3]{DC26}, where the term is used for three-colourings, $m_k(n)$ is a \emph{valid lower bound} if every $(3,k;n)$-graph has at least $m_k(n)$ edges.

For a vertex $v$ of a $(3,k;n,e)$-graph $G$, $N(v)$ is independent, so $d_v\le k-1$.
Moreover, $G_v=G-N[v]$, where $N[v]=N(v)\cup\{v\}$, is a $(3,k-1;n-1-d_v)$-graph, as an independent set of $G_v$ extends by $v$.
As no edge lies inside $N(v)$, $G_v$ has $e-Z_v$ edges, where
\[
Z_v=\sum_{u\in N(v)}d_u .
\]

As $\sum_vZ_v=\sum_vd_v^2$, valid lower bounds $m_{k-1}$ give
\[
\sum_v\big(d_v^2+m_{k-1}(n-1-d_v)\big)\le ne .
\]
This is the inequality of Graver and Yackel \cite{GrY}, on which Goedgebeur and Radziszowski \cite{GR13} base their bounds for $k\ge 11$.
It is the two-colour analogue of the density inequality \cite[Lemma~3.3]{DC26}.
As $ne=\sum_vnd_v/2$, it says
\[
\sum_v\varphi(d_v)\le 0 \qquad\text{where}\qquad \varphi(d)=d^2+m_{k-1}(n-1-d)-nd/2 .
\]

\paragraph{The bootstrap.}
We bootstrap as in the ladder of \cite[Section~5]{DC26}, with this inequality in place of the density inequality.
Level $k\ge 11$ uses the bound $B_{k-1}$ on $\R(3,k-1)$ and the valid lower bounds $m_{k-1}$ from level $k-1$.
Level $11$ starts from $B_{10}=42$ \cite[Theorem~6]{GR13} and the values and lower bounds for $e(3,10,n)$ in \cite{GR13}.

Level $k$ treats the orders $n=0,1,2,\dots$ in turn, records a valid lower bound $m_k(n)$, and stops at the first order it excludes, $B_k$, that is, the first $n$ at which it shows that no $(3,k;n)$-graph exists.
So $\R(3,k)\le B_k$.
In the same way, a value $e$, or a degree sequence, is excluded at order $n$ when no $(3,k;n,e)$-graph exists (with that degree sequence).
We run the procedure twice, giving two \emph{chains}: one from the inputs of \cite{GR13} alone, and one that also uses the lower bounds of \cite{Back3} described after the list below.
At order $n$ the procedure is as follows.
\begin{enumerate}[label=(\roman*)]
\item The admissible degrees are the $d\le\min(k-1,n-1)$ with $n-1-d\le B_{k-1}-1$.
\item An order $n\ge 1$ is excluded if $\varphi(d)>0$ at every admissible $d$ (in particular if there is none), as $\sum_v\varphi(d_v)$ would then be a sum of $n$ positive terms.
\item Otherwise let $e_0$ be the largest of the following bounds.
The first is the bound of Radziszowski and Kreher \cite[Theorem~2]{GR13}.
With $K=k-1$, it gives the exact values
\[
e(3,k,n)=\begin{cases}
0 & \text{for } n\le K,\\
n-K & \text{for } K<n\le 2K,\\
3n-5K & \text{for } 2K<n\le 5K/2,\\
5n-10K & \text{for } 5K/2<n\le 3K,\\
6n-13K & \text{for } 3K<n\le 13K/4-1,
\end{cases}
\]
and $e(3,k,n)\ge 6n-13K$ for all $n$.
The second exists when $\varphi$ is positive at the least admissible $d$ but not at every admissible $d$.
Let $(d',\varphi(d'))$ and $(d'',\varphi(d''))$, with $d'<d''$, be consecutive vertices of the lower convex hull of the points $(d,\varphi(d))$, $d$ admissible, with $\varphi(d')>0\ge\varphi(d'')$, and let $\alpha-\beta d$ be the line through them.
The second bound is $\lceil n\alpha/(2\beta)\rceil$.
This line has $\beta>0$ and lies on or below $\varphi$ at every admissible $d$, so
\[
0\ge\sum_v\varphi(d_v)\ge n\alpha-2\beta e .
\]
\item The extra step below, at the orders where the search applied it (each is recorded in the certificates), may exclude $e_0$ and some larger values of $e$.
The recorded $m_k(n)$ is the least $e\ge e_0$ that it does not exclude.
If it excludes every $e$ from $e_0$ to $\lfloor n(k-1)/2\rfloor$, the most edges possible when every degree is at most $k-1$, then $n$ is excluded.
\end{enumerate}
In the chain with \cite{Back3}, the lower bounds of \cite{Back3} on $e(3,k,n)$ for $11\le k\le 13$ and $35\le n\le 43$ are a further candidate for $e_0$ in~(iii); at $k=10$ they agree with those of \cite{GR13}.
Deleting a vertex as in \cite[Lemma~3.5(a)]{DC26} gives $m_k(n-1)+\max\{0,n-B_{k-1}\}$, as every degree is at least $n-B_{k-1}$.
This never exceeds $e_0$ in either chain, so the certificates do not use it.

\paragraph{An extra step.}
We add one step, for $e=e_0,e_0+1,\dots$ in turn.
We enumerate every degree sequence with sum $2e$ and admissible entries, each at most $e-m_k(n-1)$, that satisfies the Graver--Yackel inequality.
The cap on the entries holds as deleting $v$ leaves a $(3,k;n-1)$-graph with $e-d_v$ edges.

We then try to exclude each sequence by a small linear programme.
The \emph{class} of an entry $a$ is the set of the $c_a$ vertices of degree $a$.
The variables are the numbers $x_{ab}=x_{ba}$ of edges between the classes of $a\ne b$ and the numbers $x_{aa}$ of edges inside a class.
Counting edge ends, write $x_{a\to b}=x_{ab}$ for $a\ne b$ and $x_{a\to a}=2x_{aa}$, so that $\sum_{d_v=a}Z_v=\sum_bb\,x_{a\to b}$.
With $z_a=e-m_{k-1}(n-1-a)$, every $(3,k;n,e)$-graph with that degree sequence satisfies, for all entries $a$ and $b$:
\begin{itemize}
\item $\sum_bx_{a\to b}=a\,c_a$;
\item $\sum_bb\,x_{a\to b}\le c_az_a$, as $G_v$ has $e-Z_v\ge m_{k-1}(n-1-a)$ edges;
\item $\sum_bb\,x_{a\to b}\ge c_a\big(e-\lfloor(n-1-a)(k-2)/2\rfloor\big)$, as $G_v$ is a $(3,k-1)$-graph and so has maximum degree at most $k-2$;
\item $x_{ab}\le c_ac_b$ for $a\ne b$;
\item $\sum x_{ab}\le\lfloor N_S^2/4\rfloor$, summed over $a\le b$ in $S$, for every set $S$ of classes with $N_S$ vertices in all (Mantel's theorem);
\item $\sum_{b'\le b}x_{a\to b'}\ge t\,c_a$, with $t=t(a,b)$ as follows.
\end{itemize}

For the last constraint, let $q$ be the least entry above $b$, and let $v$ be a vertex of degree $a$.
Remove from the sequence one term $a$, that of $v$.
For $0\le j\le a$, let $L(j)$ be the sum of the $j$ smallest remaining terms at most $b$, counted with multiplicity, plus $(a-j)q$.
Put $L(j)=\infty$ if fewer than $j$ such terms remain, or if no entry of the sequence exceeds $b$ (so that $q$ is undefined) and $j<a$.
If $v$ has exactly $j$ neighbours of degree at most $b$, then $L(j)\le Z_v\le z_a$ (so no such $v$ exists when $L(j)=\infty$).
So $v$ has at least $t$ such neighbours, where $t$ is the least $j$ with $L(j)\le z_a$, and summing over the class of $a$ gives the last constraint.
If there is no such $j$, no vertex of degree $a$ exists and the sequence is excluded by this \emph{neighbour-count test}.

If every sequence is excluded, so is $e$, and the step moves on to $e+1$.
Each exclusion is certified either by the neighbour-count test or by rational Farkas multipliers, which combine these constraints, with $x\ge 0$, into an inequality $0\ge c$ with $c>0$.

\begin{theorem}[computer-assisted]\label{thm:r3k}
From the published values and lower bounds of \cite[Theorem~2 and Tables~3--5]{GR13} and $\R(3,10)\le 42$ \cite[Theorem~6]{GR13}, $\R(3,k)$ is at most the value in the row ``\cite{GR13} only'' of Table~\ref{tab:r3k} for $24\le k\le 50$.
If the lower bounds on $e(3,k,n)$ in the table of \cite[Section~5]{Back3} are also used, $\R(3,k)$ is at most the value in the row ``with \cite{Back3}''.
\end{theorem}

\begin{table}[htbp]
\centering\small
\begin{tabular}{l*{9}{c}}
\toprule
$k$ & 24 & 25 & 26 & 27 & 28 & 29 & 30 & 31 & 32\\
\midrule
\cite{GR13} only & 201 & 216 & 231 & 247 & 264 & 280 & 297 & 315 & 333\\
with \cite{Back3} & 200 & 215 & 230 & 245 & 261 & 278 & 295 & 312 & 331\\
\midrule
$k$ & 33 & 34 & 35 & 36 & 37 & 38 & 39 & 40 & 41\\
\midrule
\cite{GR13} only & 351 & 370 & 389 & 409 & 429 & 450 & 471 & 492 & 513\\
with \cite{Back3} & 349 & 369 & 389 & 408 & 429 & 450 & 471 & 492 & 513\\
\midrule
$k$ & 42 & 43 & 44 & 45 & 46 & 47 & 48 & 49 & 50\\
\midrule
both & 536 & 559 & 582 & 606 & 630 & 654 & 679 & 704 & 729\\
\bottomrule
\end{tabular}
\caption{Upper bounds for $\R(3,k)$ from Theorem~\ref{thm:r3k}.
The two rows coincide for $k=35$ and $k\ge 37$, and for $k\ge 42$ they are shown as one, ``both''.
The survey \cite{Rad26} lists no upper bounds for these $k$.
For comparison, $\R(3,k)\le\R(3,k-1)+k$, strict when both terms are even, iterated from the survey's $\R(3,23)\le 185$, gives $209$, $234$, $259$, $371$ and $1171$ for $k=24,25,26,30,50$.
Shearer's bound \cite{Sh83} $\R(3,k)\le\lfloor(k-1)/f(k-1)\rfloor+1$, with $f(d)=(d\ln d-d+1)/(d-1)^2$, gives $223$ for $k=24$ and $792$ for $k=50$.}
\label{tab:r3k}
\end{table}

\paragraph{The two chains.}
Above level $10$, each chain uses only the bounds $B_{k-1}$ that it proves itself, never the unpublished values of \cite{Rad26}.
Neither chain needs the bound $\R(3,10)\le 41$ of Angeltveit \cite{Ang25}.
The report \cite{Back3} is a preprint that gives little of its proofs, so the second row rests on inputs without published proofs.
The first row is $1$ to $3$ larger for $24\le k\le 34$ and $k=36$, and equals the second for $k=35$ and $k\ge 37$.

Without the extra step, the chain from \cite{GR13} alone gives the same bounds for $11\le k\le 50$, except $\R(3,20)\le 146$ and $\R(3,35)\le 390$.
This run is not among the certificates in the ancillary files.
The chain from \cite{GR13} alone also reaches $e(3,13,51)\ge 178$ and $e(3,14,67)\ge 335$, which the comments column of \cite[Tables~9 and~10]{GR13} credits to Backelin.
It raises seven further lower bounds of \cite[Tables~9--11]{GR13} by one, to
\begin{gather*}
e(3,13,58)\ge 261,\qquad e(3,14,69)\ge 366,\qquad e(3,14,70)\ge 382,\\
e(3,15,n)\ge 516,\ 534,\ 553\ \text{and}\ 573\quad\text{for } n=82,\dots,85.
\end{gather*}
These are values $m_k(n)$ recorded in the certificate of that chain, so the checker of Theorem~\ref{thm:r3k} verifies them.
They lie beyond the range $n\le 43$ of \cite{Back3}.
Backelin's unpublished manuscript, to which \cite{Rad26} credits $\R(3,17)\le 109$, may already contain them.

\paragraph{Checking the certificates.}
The certificates are checked by a program that does not import the generating code.
It enumerates the degree sequences itself, rebuilds each constraint from its definition, and verifies the multipliers in exact arithmetic.
\begin{itemize}
\item The chain with \cite{Back3} has one row for each level $11\le k\le 50$ and each order $0\le n<B_k$, $13{,}376$ rows in all.
Each row records $m_k(n)$ and the bound of~(iii) that gives $e_0$: the Radziszowski--Kreher bound, a value of \cite{Back3}, or the two ends of the hull edge.
At the orders where the extra step is used, it also records a certificate for every degree sequence at each excluded $e$.
A further row per level, not counted here, certifies the exclusion of $B_k$: by~(ii) at every level except $k=16$, where the extra step excludes every $e$ at $n=97$.
The chain excludes $765$ degree sequences at $53$ triples $(k,n,e)$ (including those at $k=16$, $n=97$), $686$ by the neighbour-count test and $79$ by Farkas multipliers.
\item The chain from \cite{GR13} alone has $13{,}403$ rows and excludes $766$ sequences at $50$ triples, $684$ and $82$ respectively.
All $40$ exclusions of $B_k$ in it are by~(ii).
\end{itemize}

\paragraph{Comparison with the survey.}
For $11\le k\le 23$ the two chains give the following bounds $B_k$, against the upper bounds that \cite{Rad26} lists.
\begin{center}\small\setlength{\tabcolsep}{4pt}
\begin{tabular}{l*{13}{c}}
\toprule
$k$ & 11 & 12 & 13 & 14 & 15 & 16 & 17 & 18 & 19 & 20 & 21 & 22 & 23\\
\midrule
survey \cite{Rad26} & 50 & 59 & 68 & 77 & 87 & 97 & 109 & 120 & 132 & 145 & 157 & 171 & 185\\
\cite{GR13} only & 50 & 59 & 68 & 77 & 87 & 98 & 109 & 121 & 133 & 145 & 159 & 173 & 187\\
with \cite{Back3} & 50 & 59 & 68 & 77 & 87 & 97 & 109 & 121 & 132 & 145 & 158 & 172 & 185\\
\bottomrule
\end{tabular}
\end{center}
So for $11\le k\le 23$ the procedure improves no value of \cite{Rad26}.
From \cite{GR13} alone it gives $\R(3,17)\le 109$ and $\R(3,20)\le 145$, which \cite{Rad26} credits to an unpublished manuscript of Backelin and to Lesser's thesis \cite{Les}, respectively.
With the table of \cite{Back3} it gives $\R(3,16)\le 97$ and $\R(3,19)\le 132$, the values that \cite{Rad26} lists, citing \cite{Back3}.

With the table of \cite{Back3} it also gives $\R(3,23)\le 185$.
The survey \cite{Rad26} credits this bound to a personal communication of Backelin; it now has a derivation that can be checked, given the inputs of \cite{Back3}.
At $k=18$, $21$ and $22$ the values of \cite{Rad26} come from unpublished sources, and the chain with \cite{Back3} misses each by one.

For $k=18$, both chains give $m_{18}(120)=1020=120\cdot 17/2$ from the hull.
As every degree is at most $17$, a $(3,18;120)$-graph would have exactly $1020$ edges and be $17$-regular.
Then every $G_v$ would be a $(3,17;102,731)$-graph, as $120-1-17=102$ and $1020-17^2=731$.
So $e(3,17,102)\ge 732$, one more than our $m_{17}(102)=731$, would exclude $n=120$ and give $\R(3,18)\le 120$.

\section{Lower bounds for \texorpdfstring{$\R(4,k)$}{R(4,k)}}\label{sec:r4k}

For a prime $p\equiv 1\pmod 8$ let $Q_p$ be the circulant graph on $\mathbb Z_p$ in which $x\sim y$ if and only if $x-y$ is a nonzero fourth power.
The relation is symmetric, as $-1$ is a fourth power.
Lindsay and Cain \cite[Proposition~2.1]{LinCa} computed with cliquer that $Q_{313}$ and $Q_{457}$ have clique number $3$ and independence numbers
\[
\alpha(Q_{313})=21 \qquad\text{and}\qquad \alpha(Q_{457})=24 .
\]
The ancillary checks confirm $\alpha(Q_{313})=21$ and show that $Q_{457}$ has an independent $24$-set.

Theorem~\ref{thm:r4k} does not use the upper bounds on $\alpha(Q_p)$.
An independent set of size at least $k$ contains one of size $k$.
So if a vertex set $S$ meets every independent $k$-set of $Q_p$, then $Q_p-S$ is $K_4$-free and has no independent $k$-set.
Step~3 below checks this for $(p,k)=(313,21)$ and $(457,24)$, and the $K_4$-freeness of both graphs $Q_p-S$ is also checked directly.

\begin{theorem}[witnesses checked by computer]\label{thm:r4k}
$\R(4,21)\ge 260$ and $\R(4,24)\ge 370$.
\end{theorem}

\begin{proof}
\emph{Step 1: symmetries and cosets.}
The maps $x\mapsto ax+b$ with $a$ a nonzero fourth power are automorphisms of $Q_p$.
Let $O_1,O_2,O_3$ be the cosets of the group of nonzero fourth powers in $\mathbb Z_p^*$ other than the group itself, in a fixed order, and let $r_j$ be the least element of $O_j$.

\emph{Step 2: representatives.}
Let $k\ge 2$, let $J$ be an independent $k$-set and let $u\in J$.
The translate $J-u=\{x-u:x\in J\}$ is an independent $k$-set that contains $0$, so its other elements are not fourth powers and lie in $O_1\cup O_2\cup O_3$.
Let $j$ be least with $J-u$ meeting $O_j$, say at $w$.
Then $r_j/w$ is a fourth power, and multiplying by it maps $J-u$ to an independent $k$-set that contains $0$ and $r_j$ and still avoids $O_1,\dots,O_{j-1}$, as multiplication by a fourth power fixes each coset.
Call an independent $k$-set a \emph{representative} if, for some $j$, it contains $0$ and $r_j$ and avoids $O_1,\dots,O_{j-1}$.
So every independent $k$-set is the image of a representative under one of the maps above.

\emph{Step 3: the computation.}
The sets $S$, of $54$ and $88$ vertices, were found by a weighted local search and are listed in the ancillary files.
An exact program lists all representatives for $(p,k)=(313,21)$ and $(457,24)$, by an exhaustive search for independent $(k-2)$-sets in the common non-neighbourhood of $0$ and $r_j$ outside $O_1,\dots,O_{j-1}$.
It then checks that $S$ meets every image of every representative.
So $Q_{313}-S$ is a $K_4$-free graph on $259$ vertices with no independent $21$-set, and $Q_{457}-S$ one on $369$ vertices with no independent $24$-set.
\end{proof}

\begin{proof}[Verification]
Both graphs are rebuilt from the lists of deleted vertices by programs that share no code with the search.
The $259$-vertex graph is also checked by a direct exhaustive search for $K_4$ and for independent $21$-sets, and the $369$-vertex graph by one for $K_4$ and for independent $24$-sets.
Neither search uses symmetry, so each is independent of the orbit argument above.
\end{proof}

\paragraph{Comparison with earlier bounds.}
The previous bounds are as follows.
\begin{itemize}
\item For $\R(4,21)$, the survey has $242$ \cite{SLZL}, and $\R(4,20)\ge 252$ from a recent preprint \cite{Yu26} gives $257$ by the inequality $\R(4,l)\ge\R(4,l-1)+5$ \cite[2.3(e)]{Rad26}.
\item For $\R(4,24)$, the non-neighbourhood of a vertex of $Q_{457}$ ($Q_{457}$ is $114$-regular) has $342$ vertices, is $K_4$-free and has no independent $24$-set, as such a set would extend by the vertex.
This gives $\R(4,24)\ge 343$.
The survey lists no value for $\R(4,24)$.
The inequality \cite[2.3(e)]{Rad26}, applied twice to $\R(4,22)\ge 314$ \cite{LinCa}, gives $324$.
\end{itemize}

The method does not reach $\R(4,23)\ge 320$, which would improve the bound $319$ that this inequality gives from $\R(4,22)\ge 314$.
That needs a set of at most $457-319=138$ vertices meeting every independent $23$-set of $Q_{457}$, and the smallest we found has $147$.

\section{Four colours}\label{sec:four}

For $\bt=(t_1,\dots,t_4)$ and the unit vectors $\e_i$, the four-colour form of the neighbourhood recurrence \cite[6.1(a)]{Rad26} is
\[
\R(\bt)\le\sum_{i=1}^4\R(\bt-\e_i)-2 .
\]
Indeed, consider a colouring of $K_n$ with no colour-$j$ $K_{t_j}$ for any $j$.
The colour-$i$ neighbourhood $N_i(v)$ of a vertex $v$, whose size is the colour-$i$ degree of $v$, contains no colour-$i$ $K_{t_i-1}$, which would extend by $v$, and no colour-$j$ $K_{t_j}$ for $j\ne i$.
So it has at most $\R(\bt-\e_i)-1$ vertices.
Summing over $i$ gives
\[
n-1\le\sum_i(\R(\bt-\e_i)-1) .
\]
So such a colouring has $n\le\sum_i\R(\bt-\e_i)-3$, which is the recurrence.

The recurrence improves to $\R(\bt)\le\sum_i\R(\bt-\e_i)-3$ when $\sum_i\R(\bt-\e_i)$ is even and some $\R(\bt-\e_i)$ is even.
Indeed, take a colouring as above of $K_n$ with $n=\sum_i\R(\bt-\e_i)-3$, which is odd.
Then the display holds with equality, so every vertex has colour-$i$ degree exactly $\R(\bt-\e_i)-1$ for each $i$.
For an $i$ with $\R(\bt-\e_i)$ even this degree is odd, and a graph on an odd number of vertices cannot have all its degrees odd.

Both statements hold with the same proofs when each $\R(\bt-\e_i)$ is replaced by a proven upper bound for it.
An entry equal to $2$ may be dropped, leaving a number with fewer colours, as a colouring with no colour-$i$ $K_2$ has no colour-$i$ edges.

\begin{theorem}[from computer-assisted inputs]\label{thm:four}
$\R(3,3,4,4)\le 442$, $\R(3,4,4,4)\le 1521$ and $\R(4,4,4,4)\le 6082$.
\end{theorem}

\begin{proof}
With $\R(3,4,4)\le 73$ \cite[Theorem~8.3]{DC26} and $\R(3,3,3,4)\le 149$ \cite{Rad26}, the recurrence gives
\[
\R(3,3,4,4)\le 2\cdot 73+2\cdot 149-2=442 .
\]
Parity does not apply, as every term is odd.

With $\R(4,4,4)\le 198$ \cite[Theorem~8.2]{DC26} it then gives
\[
\R(3,4,4,4)\le 198+3\cdot 442-2=1522 .
\]
The sum $198+3\cdot 442=1524$ is even, as are all four terms, so parity gives $\R(3,4,4,4)\le 1524-3=1521$.

Finally,
\[
\R(4,4,4,4)\le 4\cdot 1521-2=6082 ,
\]
and parity does not apply, as $1521$ is odd.
\end{proof}

\paragraph{The inputs and earlier bounds.}
The two inputs from \cite{DC26} are computer-assisted.
The survey's $\R(3,3,3,4)\le 149$ is itself the recurrence with parity, from $\R(3,3,4)=30$ \cite{CFIM16} and $\R(3,3,3,3)\le 62$ \cite{FKR04}: as $3\cdot 30+62=152$ and $30$ are even, it gives $152-3=149$.
From the survey's $\R(3,4,4,4)\le 1576$, the recurrence with parity gives $\R(4,4,4,4)\le 4\cdot 1576-3=6301$, as $4\cdot 1576$ and all four terms are even.
The survey lists $6306$ and notes that $6301$ is implicit in \cite{Boz26}.

With the survey's $\R(3,3,4,4)\le 450$ in place of $442$, that is, without $\R(3,4,4)\le 73$, the bound $\R(4,4,4)\le 198$ alone gives $\R(3,4,4,4)\le 1545$ and $\R(4,4,4,4)\le 6178$.
The survey's $\R(3,4,4,4)\le 1576$, which it marks as implicit in \cite{Boz26}, comes from Boza's analogue of the parity clause \cite[Theorem~2.1]{Boz26}, which counts the triangles of one colour.
It lowers the right-hand side $P=\sum_i\R(\bt-\e_i)-2$ of the recurrence by one when three conditions modulo~$3$ all hold, the first being $P\not\equiv 1\pmod 3$.
From Boza's $\R(4,4,4)\le 229$ and the survey's $\R(3,3,4,4)\le 450$ it gives $229+3\cdot 450-2-1=1576$.
With our bounds it gives nothing more.
At the three steps $P$ is $442$, $1522$ and $6082$ (before any parity improvement), all $\equiv 1\pmod 3$, so the first condition fails.

\paragraph{Four triangle-free colours.}
We also tried to improve $\R(3,3,3,3)\le 62$ \cite{FKR04}, that is, to exclude colourings of $K_{61}$ with four triangle-free colours, by certificates analogous to those for $\R(4,4,4)$ in \cite[Sections~7 and~8]{DC26}.
Such a certificate exists only if a flag relaxation, here with four- or five-vertex patterns and with four-colour analogues of the local caps of \cite[Lemmas~2.1 and~2.3]{DC26} among its constraints, is infeasible.
As in Razborov's flag algebras \cite{Raz07}, the unknowns of the relaxation are the densities of the colourings of $K_4$ or $K_5$, the patterns, in a putative colouring of $K_{61}$.
They are subject to linear constraints from the caps, such as $|N_i(v)\cap N_j(u)|\le\R(3,3)-1=5$ for $u\in N_i(v)$ and $j\ne i$ (an edge of colour $i$ or $j$ inside this set would form a triangle of that colour with $v$ or with $u$, so the set carries only the two other colours, both triangle-free), and to positive-semidefiniteness constraints.
A certificate is a dual solution that proves the relaxation infeasible.
For comparison, the six-vertex computation of Lidick\'y and Pfender \cite[Appendix~A]{LP21}, whose relaxation has no caps, gives only $\R(3,3,3,3)\le 65$.

At $n=61$ each colour-$i$ degree is at most $\R(3,3,3)-1=16$ \cite{GG55}, and the four colour degrees of a vertex sum to $60$.
So the \emph{profile} of a vertex, its sorted quadruple of colour degrees, is one of
\[
(12,16,16,16),\quad (13,15,16,16),\quad (14,14,16,16),\quad (14,15,15,16),\quad (15,15,15,15).
\]

Both relaxations are numerically feasible.
Floating-point solvers found an assignment of pattern densities meeting all their constraints.
A case split by vertex profiles does not help either.
The solvers found such an assignment even with the added requirement that every vertex has profile $(12,16,16,16)$, so the case of any such split that allows this profile at every vertex is not excluded.
The four-vertex relaxation was solved as a linear programme by HiGHS \cite{HH18}, with its semidefinite constraint imposed by eigenvector cuts until no eigenvalue was below $-10^{-8}$.
The five-vertex relaxation was solved as a semidefinite programme by Clarabel \cite{GC24} at tolerance $10^{-8}$.
So we obtained no such certificate from these relaxations, and no bound follows.

\section{Verification, literature and limitations}\label{sec:verify}

\paragraph{Ancillary files.}
This paper and \cite{DC26} have the same ancillary files, which are also in the directory \texttt{ancillary/} of the repository
\begin{center}\url{https://github.com/Lachy-Dauth/ramsey-neighbourhood-counting}.\end{center}
They contain the programs and data, among them:
\begin{itemize}
\item \texttt{two\_colour\_r3k/} for Theorem~\ref{thm:r3k}, with the inputs from \cite{GR13} and \cite{Back3} and the certificates of both chains;
\item \texttt{two\_colour\_r4k/} for Theorem~\ref{thm:r4k}, with the lists of deleted vertices, \texttt{build\_power\_sub.py}, which rebuilds the two graphs from them, the exact checkers \texttt{rcheck.c} (direct search) and \texttt{qpsym.c} (orbit argument), and \texttt{verify\_all.sh}, which checks the SHA-256 hashes of the two graphs and runs the checkers;
\item \texttt{flag5\_r444/} and \texttt{flag5\_r34k/}, which check the inputs $\R(4,4,4)\le 198$ and $\R(3,4,4)\le 73$ of Theorem~\ref{thm:four} from \cite{DC26}.
\end{itemize}
The command \texttt{sh run\_all.sh} runs the quick checks of both papers in a few minutes.
The command \texttt{sh run\_all.sh --full} adds the long ones, among them the checks of these two inputs and the direct search of the $369$-vertex graph for independent $24$-sets.

The checkers for this paper, in \texttt{two\_colour\_r3k/} and \texttt{two\_colour\_r4k/}, are written in Python and C with their standard libraries, and the C programs use POSIX threads.
To run only these checks, use the commands in \texttt{two\_colour\_r3k/README.md} and the script \texttt{two\_colour\_r4k/verify\_all.sh}.
The script \texttt{run\_all.sh} also needs a C++17 compiler, for the checks of \cite{DC26}, and uses NumPy, when it is installed, for an optional second checker in \texttt{flag5\_r444/} and its negative controls \cite[Section~10]{DC26}.
Negative controls are all rejected: a weakened input, a zeroed multiplier, a removed sequence certificate, a value raised by one, a claimed bound below the proved one such as $\R(3,24)\le 199$, and a deleted vertex put back.

\paragraph{Literature.}
We searched arXiv and the web up to 23 October 2026 for earlier statements of our new bounds, and found none.
\begin{itemize}
\item For $\R(3,k)$ with $k\ge 24$ the survey lists no upper bounds, and the upper bounds of \cite{GR13} and of Lesser's thesis \cite{Les} stop at $k=16$ and $k=21$.
Backelin's unpublished manuscript (Section~\ref{sec:r3k}) may contain bounds of this kind.
\item For $\R(4,21)$ the best previous bound we know of, $257$, follows from $\R(4,20)\ge 252$ in the preprint \cite{Yu26}, which we have not checked.
Without it the bound is $242$ \cite{SLZL}.
\item For $\R(4,24)$ the survey lists nothing, and $343$ is implicit in \cite{LinCa}.
\item For the four-colour numbers the previous bounds are those of \cite{Rad26} and \cite{Boz26} in Table~\ref{tab:results}.
\end{itemize}
Such a search cannot establish priority, and we would welcome pointers to any earlier source.

\paragraph{Limitations.}
For $11\le k\le 23$ the bootstrap of Section~\ref{sec:r3k} improves none of the bounds in \cite{Rad26}.
At $k=18$ it would give $\R(3,18)\le 120$ if its own lower bound $e(3,17,102)\ge 731$ were one larger.
The second row of Table~\ref{tab:r3k} rests on the lower bounds of \cite{Back3}, whose proofs are largely unpublished.

The deletion method of Section~\ref{sec:r4k} does not reach $\R(4,23)\ge 320$.
The four-colour bounds rest on $\R(3,4,4)\le 73$ and $\R(4,4,4)\le 198$ \cite[Theorems~8.3 and~8.2]{DC26}, which are computer-assisted, and on $\R(3,3,3,4)\le 149$, which \cite[6.1(a)]{Rad26} derives from $\R(3,3,4)=30$ \cite{CFIM16} and $\R(3,3,3,3)\le 62$ \cite{FKR04}, both computer-assisted.

\section*{Acknowledgements}
\pdfbookmark[1]{Acknowledgements}{acknowledgements}
Much of the work reported here was carried out with two AI systems, used under the direction of the author: OpenAI's Codex, based on GPT-6, and Claude Opus~5.5, a large language model developed by Anthropic.
They searched for the arguments and certificates, wrote the search programs for this paper (not included) and the checking programs in the ancillary files, and drafted the text.
Every certificate and witness is checked in exact arithmetic by the programs in the ancillary files, which can be run without either AI system.
The author takes full responsibility for the contents of this paper.

\pdfbookmark[1]{References}{references}
\begin{thebibliography}{99}\small
\bibitem{Ang25} V.~Angeltveit, $R(3,10)\le 41$, \emph{Electron.\ J.\ Combin.} 32(4) (2025), \#P4.30.
\bibitem{Back3} J.~Backelin, Edge number report 1: state of the art estimates for $n\le 43$, preprint, arXiv:1410.1843, 2014.
\bibitem{Boz26} L.~Boza, New upper bounds for the classical Ramsey numbers $R(4,4,4)$, $R(3,4,5)$ and $R(3,3,6)$, arXiv:2603.10851, 2026.
\bibitem{CFIM16} M.~Codish, M.~Frank, A.~Itzhakov and A.~Miller, Computing the Ramsey number $R(4,3,3)$ using abstraction and symmetry breaking, \emph{Constraints} 21 (2016), 375--393.
\bibitem{DC26} L.~Dauth, Neighbourhood counting and new bounds for small three-colour Ramsey numbers, companion paper, 2026. The theorem and section numbers cited are those of its first arXiv version.
\bibitem{FKR04} S.~E.~Fettes, R.~L.~Kramer and S.~P.~Radziszowski, An upper bound of 62 on the classical Ramsey number $R(3,3,3,3)$, \emph{Ars Combin.} 72 (2004), 41--63.
\bibitem{GR13} J.~Goedgebeur and S.~P.~Radziszowski, New computational upper bounds for Ramsey numbers $R(3,k)$, \emph{Electron.\ J.\ Combin.} 20(1) (2013), \#P30.
\bibitem{GC24} P.~J.~Goulart and Y.~Chen, Clarabel: an interior-point solver for conic programs with quadratic objectives, \emph{Math.\ Program.\ Comput.} (2026), doi:10.1007/s12532-026-00320-7; arXiv:2405.12762.
\bibitem{GrY} J.~E.~Graver and J.~Yackel, Some graph theoretic results associated with Ramsey's theorem, \emph{J.\ Combin.\ Theory} 4 (1968), 125--175.
\bibitem{GG55} R.~E.~Greenwood and A.~M.~Gleason, Combinatorial relations and chromatic graphs, \emph{Canad.\ J.\ Math.} 7 (1955), 1--7.
\bibitem{HH18} Q.~Huangfu and J.~A.~J.~Hall, Parallelizing the dual revised simplex method, \emph{Math.\ Program.\ Comput.} 10 (2018), 119--142.
\bibitem{Les} A.~S.~Lesser, Theoretical and computational aspects of Ramsey theory, Examensarbeten i Matematik 3, Stockholms Universitet, 2001, \url{https://www2.math.su.se/ramseytheory/sciramseyams.pdf}.
\bibitem{LP21} B.~Lidick\'y and F.~Pfender, Semidefinite programming and Ramsey numbers, \emph{SIAM J.\ Discrete Math.} 35 (2021), 2328--2344; arXiv:1704.03592.
\bibitem{LinCa} M.~Lindsay and J.~W.~Cain, Improved lower bounds on the classical Ramsey numbers $R(4,22)$ and $R(4,25)$, preprint, arXiv:1510.06102, 2015.
\bibitem{Rad26} S.~P.~Radziszowski, Small Ramsey numbers, \emph{Electron.\ J.\ Combin.}, Dynamic Survey DS1.18 (revision 18, 24 April 2026), doi:10.37236/21.
\bibitem{Raz07} A.~A.~Razborov, Flag algebras, \emph{J.\ Symbolic Logic} 72 (2007), 1239--1282.
\bibitem{Sh83} J.~B.~Shearer, A note on the independence number of triangle-free graphs, \emph{Discrete Math.} 46 (1983), 83--87.
\bibitem{SLZL} W.~Su, H.~Luo, Z.~Zhang and G.~Li, New lower bounds of fifteen classical Ramsey numbers, \emph{Australas.\ J.\ Combin.} 19 (1999), 91--99.
\bibitem{Yu26} C.~Yu, Ramsey number $R(4,20)\ge 252$, preprint, arXiv:2608.18169, 2026.
\end{thebibliography}

\end{document}
